% OWNER: theory-b
% STATUS: outline
% SOURCES: notes/SPIS_TRESCI_ROBOCZY.md § Rozdz. 4
\chapter{Osobowość a preferencje przestrzenne}
\label{ch04:top}
\ChapterStatus{outline}{theory-b}

\section{Model Wielkiej Piątki}
\label{ch04:sec:bigfive}
\Todo{Konstrukt, stabilność, uzasadnienie wyboru modelu.}
\NeedsCite{john2008bigfive}

\section{IPIP i IPIP-NEO-120}
\label{ch04:sec:ipip}
\Todo{Facety jako poziom szczegółowości; implementacja w IDA.}

\section{Literatura: osobowość a preferencje środowiskowe}
\label{ch04:sec:literatura}
\Todo{Przegląd korelacji osobowość--estetyka/środowisko.}

\section{Mapowanie facetów na parametry wnętrz}
\label{ch04:sec:mapowanie}
\Todo{Kolory, złożoność, funkcja, harmonia — odwołanie do załącznika~\ref{app:c:top}.}

\section{Krytyka i ograniczenia}
\label{ch04:sec:krytyka}
\Todo{Korelacje vs przyczynowość; ryzyko determinizmu.}

\section{Uzasadnienie włączenia osobowości do pipeline IDA}
\label{ch04:sec:uzasadnienie}
\Todo{Źródło \texttt{personality} w macierzy; związek z pytaniami badawczymi dot.\ osobowości.}
